\documentclass[a4paper,12pt]{article}
\usepackage[utf8]{inputenc}
\usepackage[T1]{fontenc}
\usepackage[spanish,es-nodecimaldot]{babel}
\usepackage{geometry}
\geometry{margin=2.2cm}
\usepackage{amsmath,amssymb}
\usepackage{siunitx}
\usepackage{booktabs}
\usepackage{float}
\usepackage{array}
\usepackage{graphicx}
\usepackage{tikz}
\usetikzlibrary{calc}
\sisetup{output-decimal-marker={,},per-mode=symbol}

\newcommand{\triangulopot}[4]{%
\begin{tikzpicture}[x=0.042cm,y=0.042cm,baseline=(current bounding box.center)]
\def\P{#1}
\def\Q{#2}
\coordinate (O) at (0,0);
\coordinate (A) at (\P,0);
\coordinate (B) at (\P,\Q);
\draw[->] (-10,0) -- (\P+18,0) node[right] {$P$};
\ifdim #2pt > 0pt
  \draw[->] (\P,0) -- (\P,\Q+18) node[above] {$Q$};
\fi
\draw[thick] (O) -- (A) -- (B) -- cycle;
\node[below] at ($(O)!0.5!(A)$) {\footnotesize $#1\,\si{W}$};
\ifdim #2pt > 0pt
  \node[right] at ($(A)!0.5!(B)$) {\footnotesize $#2\,\si{VAr}$};
  \node[above left] at ($(O)!0.5!(B)$) {\footnotesize $#3\,\si{VA}$};
\else
  \node[above] at ($(O)!0.5!(A)$) {\footnotesize $#3\,\si{VA}$};
\fi
\node[below=8pt] at ($(O)!0.5!(A)+(0,-10)$) {\footnotesize #4};
\end{tikzpicture}%
}

\newcommand{\bloquesimple}[3]{%
\begin{tikzpicture}[scale=0.95,every node/.style={font=\small}]
\draw (0,3.5) -- (6.2,3.5);
\draw (0,0) -- (6.2,0);
\draw (0,3.5) -- (0,2.2);
\draw (0,1.3) -- (0,0);
\draw (0,1.75) circle (0.45);
\node at (0,1.75) {$V_s$};
\draw (3,3.5) -- (3,2.25);
\draw (3,1.25) -- (3,0);
\draw (2.0,2.25) rectangle (4.0,1.25);
\node[align=center] at (3,1.75) {#1\\#2};
\node at (3.1,-0.6) {#3};
\end{tikzpicture}%
}

\newcommand{\bloquedoble}[4]{%
\begin{tikzpicture}[scale=0.95,every node/.style={font=\small}]
\draw (0,3.5) -- (6.8,3.5);
\draw (0,0) -- (6.8,0);
\draw (0,3.5) -- (0,2.2);
\draw (0,1.3) -- (0,0);
\draw (0,1.75) circle (0.45);
\node at (0,1.75) {$V_s$};
\draw (2.0,3.5) -- (2.0,2.25);
\draw (2.0,1.25) -- (2.0,0);
\draw (1.1,2.25) rectangle (2.9,1.25);
\node[align=center] at (2.0,1.75) {#1};
\draw (5.0,3.5) -- (5.0,2.25);
\draw (5.0,1.25) -- (5.0,0);
\draw (3.7,2.25) rectangle (6.3,1.25);
\node[align=center] at (5.0,1.75) {#2\\#3};
\node at (3.8,-0.6) {#4};
\end{tikzpicture}%
}

\newcommand{\bloquetriple}[5]{%
\begin{tikzpicture}[scale=0.95,every node/.style={font=\small}]
\draw (0,3.5) -- (8.7,3.5);
\draw (0,0) -- (8.7,0);
\draw (0,3.5) -- (0,2.2);
\draw (0,1.3) -- (0,0);
\draw (0,1.75) circle (0.45);
\node at (0,1.75) {$V_s$};
\draw (1.8,3.5) -- (1.8,2.25);
\draw (1.8,1.25) -- (1.8,0);
\draw (1.0,2.25) rectangle (2.6,1.25);
\node[align=center] at (1.8,1.75) {#1};
\draw (4.6,3.5) -- (4.6,2.25);
\draw (4.6,1.25) -- (4.6,0);
\draw (3.3,2.25) rectangle (5.9,1.25);
\node[align=center] at (4.6,1.75) {#2\\#3};
\draw (7.2,3.5) -- (7.2,2.25);
\draw (7.2,1.25) -- (7.2,0);
\draw (6.3,2.25) rectangle (8.1,1.25);
\node[align=center] at (7.2,1.75) {#4};
\node at (4.8,-0.6) {#5};
\end{tikzpicture}%
}

\begin{document}

\section{Desarrollo}

De acuerdo con la guía de elaboración, esta sección reúne la descripción de lo realizado en el laboratorio, los esquemáticos de las configuraciones implementadas, los resultados experimentales y las magnitudes derivadas necesarias para el posterior análisis. Para el TP 3 se tomaron, en cada configuración, las mediciones de tensión \(V\), corriente \(I\) y potencia activa \(P\), y a partir de ellas se obtuvieron la potencia aparente \(S\), la potencia reactiva \(Q\), el factor de potencia \(\cos\varphi\) y la impedancia equivalente de carga \(Z_L\), tal como indica la consigna del trabajo práctico.

Las expresiones utilizadas para las magnitudes indirectas fueron:
\[
S = VI,\qquad
\cos\varphi = \frac{P}{S},\qquad
Q = \sqrt{S^2-P^2},\qquad
Z_L = \frac{V}{I}.
\]
En aquellos casos en los que la lectura instrumental produjo un valor \(P>VI\), se adoptó para la presentación \(\cos\varphi \approx 1\), \(Q\approx 0\) y \(S\approx P\), interpretando la discrepancia como consecuencia de la resolución finita de los instrumentos analógicos.

\subsection{Parámetros utilizados}

\begin{table}[H]
    \centering
    \caption{Parámetros medidos y valores utilizados en el desarrollo experimental.}
    \label{tab:parametros_tp3}
    \begin{tabular}{>{\raggedright\arraybackslash}p{0.58\textwidth}c}
        \toprule
        Magnitud & Valor \\
        \midrule
        Tensión de alimentación de los ensayos & \SI{100}{\volt} \\
        Frecuencia de la red & \SI{50}{\hertz} \\
        Resistencia equivalente del banco de lámparas & \SI{166,7}{\ohm} \\
        Resistencia serie de la bobina L1 & \SI{65,4}{\ohm} \\
        Resistencia serie de la bobina L2 & \SI{22,1}{\ohm} \\
        Capacidad teórica para compensar el caso base & \SI{7,05}{\micro\farad} \\
        Capacidad real implementada con el banco & \SI{8,0}{\micro\farad} \\
        Inductancia estimada de L1 (medición complementaria) & \SI{3,480}{\henry} \\
        Inductancia estimada de L2 (medición complementaria) & \SI{0,338}{\henry} \\
        \bottomrule
    \end{tabular}
\end{table}

\subsection{Paso 1: rama resistiva}

\begin{figure}[H]
    \centering
    \bloquesimple{$R$}{\SI{166,7}{\ohm}}{Solo $S_R$ cerrada}
    \caption{Configuración del paso 1: únicamente se conecta la rama resistiva equivalente al banco de lámparas.}
    \label{fig:tp3_paso1}
\end{figure}

\begin{table}[H]
    \centering
    \caption{Resultados del paso 1.}
    \label{tab:tp3_paso1}
    \begin{tabular}{>{\raggedright\arraybackslash}p{0.33\textwidth}ccccccc}
        \toprule
        Configuración & $V$ [V] & $I$ [A] & $P$ [W] & $S$ [VA] & $Q$ [VAr] & $\cos\varphi$ & $Z_L$ [$\Omega$] \\
        \midrule
        Solo rama resistiva ($S_R$ cerrada) & 100 & 0,60 & 62 & 62 & $\approx 0$ & $\approx 1,00$ & 166,7 \\
        \bottomrule
    \end{tabular}
\end{table}

\subsection{Paso 2: rama resistiva + bobina con núcleo de aire}

\begin{figure}[H]
    \centering
    \bloquedoble{$R$}{$R_L$}{$L$}{$S_R$ y $S_L$ cerradas}
    \caption{Configuración general de la rama resistivo--inductiva empleada en los pasos 2, 4 y 5.}
    \label{fig:tp3_rl_general}
\end{figure}

En el paso 2 se utilizó la configuración de la Figura~\ref{fig:tp3_rl_general}, con $R_L=\SI{65,4}{\ohm}$ y la bobina $L_1$ con núcleo de aire.

\begin{table}[H]
    \centering
    \caption{Resultados del paso 2.}
    \label{tab:tp3_paso2}
    \begin{tabular}{>{\raggedright\arraybackslash}p{0.33\textwidth}ccccccc}
        \toprule
        Configuración & $V$ [V] & $I$ [A] & $P$ [W] & $S$ [VA] & $Q$ [VAr] & $\cos\varphi$ & $Z_L$ [$\Omega$] \\
        \midrule
        Rama resistiva + L1 (núcleo de aire) & 100 & 0,70 & 66,4 & 70 & 22,2 & 0,949 & 142,9 \\
        \bottomrule
    \end{tabular}
\end{table}

A partir de esta configuración se calculó la capacidad necesaria para compensar el factor de potencia. Con \(V=\SI{100}{\volt}\), \(I=\SI{0,70}{\ampere}\), \(P=\SI{66,4}{\watt}\) y \(Q\approx\SI{22,16}{\volt\ampere reactive}\), se obtuvo:
\[
C_{teo}=\frac{Q}{\omega V^2}=\frac{22,16}{2\pi\cdot 50\cdot 100^2}\approx \SI{7,05}{\micro\farad}.
\]
El valor implementado con el banco fue \(\,C_{real}=\SI{8,0}{\micro\farad}\).

\subsection{Paso 3: compensación capacitiva}

\begin{figure}[H]
    \centering
    \bloquetriple{$R$}{$R_L$}{$L$}{$C$}{$S_R$, $S_L$ y $S_C$ cerradas}
    \caption{Configuración general de la rama resistivo--inductiva compensada con el banco capacitivo.}
    \label{fig:tp3_rlc_general}
\end{figure}

En el paso 3 se empleó la configuración de la Figura~\ref{fig:tp3_rlc_general}, tomando $R_L=\SI{65,4}{\ohm}$, la bobina $L_1$ y la capacidad implementada con el banco. La fila de referencia de la Tabla~\ref{tab:tp3_paso3} corresponde al mismo circuito, pero asignando $C=C_{teo}=\SI{7,05}{\micro\farad}$ en lugar de $C=C_{real}=\SI{8,0}{\micro\farad}$.

\begin{table}[H]
    \centering
    \caption{Resultados del paso 3, comparando la capacidad real implementada y la capacidad teórica de referencia.}
    \label{tab:tp3_paso3}
    \begin{tabular}{>{\raggedright\arraybackslash}p{0.36\textwidth}ccccccc}
        \toprule
        Configuración & $V$ [V] & $I$ [A] & $P$ [W] & $S$ [VA] & $Q$ [VAr] & $\cos\varphi$ & $Z_L$ [$\Omega$] \\
        \midrule
        Compensación con $C_{real}=\SI{8,0}{\micro\farad}$ & 100 & 0,657 & 68 & 68 & $\approx 0$ & $\approx 1,00$ & 152,2 \\
        Referencia con $C_{teo}=\SI{7,05}{\micro\farad}$ & 100 & 0,66 & 68 & 68 & $\approx 0$ & $\approx 1,00$ & 151,5 \\
        \bottomrule
    \end{tabular}
\end{table}

\subsection{Paso 4: inserción de núcleo de hierro}

Para el paso 4 se reutilizó la configuración general de la Figura~\ref{fig:tp3_rl_general}, manteniendo la rama resistiva y reemplazando la rama inductiva por la correspondiente a cada ensayo con núcleo de hierro. En particular, para el primer caso se tomó $R_L=\SI{65,4}{\ohm}$ con la bobina $L_1$, mientras que para el segundo se tomó $R_L=\SI{22,1}{\ohm}$ con la bobina $L_2$.

\begin{table}[H]
    \centering
    \caption{Resultados del paso 4 con núcleo de hierro.}
    \label{tab:tp3_paso4}
    \begin{tabular}{>{\raggedright\arraybackslash}p{0.32\textwidth}ccccccc}
        \toprule
        Configuración & $V$ [V] & $I$ [A] & $P$ [W] & $S$ [VA] & $Q$ [VAr] & $\cos\varphi$ & $Z_L$ [$\Omega$] \\
        \midrule
        Núcleo de hierro con L1 & 100 & 0,60 & 64 & 64 & $\approx 0$ & $\approx 1,00$ & 166,7 \\
        Núcleo de hierro con L2 & 100 & 0,72 & 68 & 72 & 23,7 & 0,944 & 138,9 \\
        \bottomrule
    \end{tabular}
\end{table}

\subsection{Paso 5: inserción de núcleo de aluminio}

Para el paso 5 se volvió a utilizar la configuración general de la Figura~\ref{fig:tp3_rl_general}, modificando únicamente la rama inductiva para insertar un núcleo de aluminio. En los ensayos se consideraron, respectivamente, $R_L=\SI{65,4}{\ohm}$ con la bobina $L_1$ y $R_L=\SI{22,1}{\ohm}$ con la bobina $L_2$.

\begin{table}[H]
    \centering
    \caption{Resultados del paso 5 con núcleo de aluminio.}
    \label{tab:tp3_paso5}
    \begin{tabular}{>{\raggedright\arraybackslash}p{0.32\textwidth}ccccccc}
        \toprule
        Configuración & $V$ [V] & $I$ [A] & $P$ [W] & $S$ [VA] & $Q$ [VAr] & $\cos\varphi$ & $Z_L$ [$\Omega$] \\
        \midrule
        Núcleo de aluminio con L1 & 100 & 0,76 & 72 & 76 & 24,3 & 0,947 & 131,6 \\
        Núcleo de aluminio con L2 & 100 & 1,262 & 91 & 126,2 & 87,4 & 0,721 & 79,2 \\
        \bottomrule
    \end{tabular}
\end{table}

\subsection{Medición complementaria: inductor puro}

Además de las cinco configuraciones pedidas en la consigna, se realizó una medición complementaria conectando únicamente el inductor para estimar su inductancia a partir de la corriente, la potencia y la resistencia serie medida.

\begin{figure}[H]
    \centering
    \bloquesimple{$R_{L1}+L_1$}{Solo rama inductiva}{Solo $S_L$ cerrada}
    \caption{Esquema de la medición complementaria para el inductor puro.}
    \label{fig:tp3_inductor_puro}
\end{figure}

\begin{table}[H]
    \centering
    \caption{Resultados de la medición complementaria del inductor puro.}
    \label{tab:tp3_inductor_puro}
    \begin{tabular}{>{\raggedright\arraybackslash}p{0.25\textwidth}ccccccccc}
        \toprule
        Configuración & $V$ [V] & $I$ [A] & $P$ [W] & $S$ [VA] & $Q$ [VAr] & $\cos\varphi$ & $Z_L$ [$\Omega$] & $R_s$ [$\Omega$] & $L$ [H] \\
        \midrule
        Inductor puro L1 & 100 & 0,09 & 1,6 & 9 & 8,9 & 0,178 & 1111,1 & 65,4 & 3,480 \\
        Inductor puro L2 & 100 & 0,92 & 20 & 92 & 89,8 & 0,217 & 108,7 & 22,1 & 0,338 \\
        \bottomrule
    \end{tabular}
\end{table}

\subsection{Tablas resumen}

\begin{table}[H]
    \centering
    \caption{Resumen de las mediciones directas de tensión, corriente y potencia para todas las configuraciones principales del TP 3.}
    \label{tab:tp3_resumen_directas}
    \begin{tabular}{>{\raggedright\arraybackslash}p{0.38\textwidth}ccc}
        \toprule
        Configuración & $V$ [V] & $I$ [A] & $P$ [W] \\
        \midrule
        Solo rama resistiva ($S_R$) & 100 & 0,60 & 62 \\
        Rama resistiva + L1 (aire) & 100 & 0,70 & 66,4 \\
        Compensación con $C_{real}=\SI{8,0}{\micro\farad}$ & 100 & 0,657 & 68 \\
        Referencia con $C_{teo}=\SI{7,05}{\micro\farad}$ & 100 & 0,66 & 68 \\
        Núcleo de hierro con L1 & 100 & 0,60 & 64 \\
        Núcleo de hierro con L2 & 100 & 0,72 & 68 \\
        Núcleo de aluminio con L1 & 100 & 0,76 & 72 \\
        Núcleo de aluminio con L2 & 100 & 1,262 & 91 \\
        \bottomrule
    \end{tabular}
\end{table}

\begin{table}[H]
    \centering
    \caption{Resumen de magnitudes derivadas para las configuraciones principales del TP 3.}
    \label{tab:tp3_resumen_indirectas}
    \begin{tabular}{>{\raggedright\arraybackslash}p{0.38\textwidth}cccc}
        \toprule
        Configuración & $Q$ [VAr] & $S$ [VA] & $\cos\varphi$ & $Z_L$ [$\Omega$] \\
        \midrule
        Solo rama resistiva ($S_R$) & $\approx 0$ & 62 & $\approx 1,00$ & 166,7 \\
        Rama resistiva + L1 (aire) & 22,2 & 70 & 0,949 & 142,9 \\
        Compensación con $C_{real}=\SI{8,0}{\micro\farad}$ & $\approx 0$ & 68 & $\approx 1,00$ & 152,2 \\
        Referencia con $C_{teo}=\SI{7,05}{\micro\farad}$ & $\approx 0$ & 68 & $\approx 1,00$ & 151,5 \\
        Núcleo de hierro con L1 & $\approx 0$ & 64 & $\approx 1,00$ & 166,7 \\
        Núcleo de hierro con L2 & 23,7 & 72 & 0,944 & 138,9 \\
        Núcleo de aluminio con L1 & 24,3 & 76 & 0,947 & 131,6 \\
        Núcleo de aluminio con L2 & 87,4 & 126,2 & 0,721 & 79,2 \\
        \bottomrule
    \end{tabular}
\end{table}

\subsection{Triángulos de potencia}

\begin{figure}[H]
    \centering
    \begin{minipage}{0.31\textwidth}\centering\triangulopot{62}{0}{62}{Solo $R$}\end{minipage}
    \begin{minipage}{0.31\textwidth}\centering\triangulopot{66.4}{22.2}{70}{$R\parallel L_1$}\end{minipage}
    \begin{minipage}{0.31\textwidth}\centering\triangulopot{68}{0}{68}{$R\parallel L_1\parallel C$}\end{minipage}

    \vspace{0.35cm}

    \begin{minipage}{0.31\textwidth}\centering\triangulopot{64}{0}{64}{Hierro -- L1}\end{minipage}
    \begin{minipage}{0.31\textwidth}\centering\triangulopot{68}{23.7}{72}{Hierro -- L2}\end{minipage}
    \begin{minipage}{0.31\textwidth}\centering\triangulopot{72}{24.3}{76}{Aluminio -- L1}\end{minipage}

    \vspace{0.35cm}

    \begin{minipage}{0.31\textwidth}\centering\triangulopot{91}{87.4}{126.2}{Aluminio -- L2}\end{minipage}
    \caption{Triángulos de potencia construidos a partir de los valores obtenidos en cada configuración principal.}
    \label{fig:tp3_triangulos}
\end{figure}

\end{document}
